\documentclass[aps,prl,twocolumn,superscriptaddress,nofootinbib]{revtex4-2}

\usepackage{amsmath,amssymb}
\usepackage{hyperref}
\usepackage{booktabs}
\usepackage{amsthm}
\newtheorem{theorem}{Theorem}
\newtheorem{observation}[theorem]{Observation}

\newcommand{\Z}{\mathbb{Z}}
\newcommand{\Q}{\mathbb{Q}}
\newcommand{\Spec}{\mathrm{Spec}}
\newcommand{\cd}{\mathrm{cd}}

\begin{document}

\title{The Proton-Electron Mass Ratio as a Gaussian Norm:\\
$m_p/m_e = 12 \times 153 = 1836$ from the $j$-Invariant\\
and the Cohomological Dimension of $\Spec(\Z)$}

\author{Wright Brothers}
\affiliation{Independent Research}
\date{\today}

\begin{abstract}
We observe that the proton-to-electron mass ratio is given,
to a precision of 0.008\%, by
\begin{equation*}
  \frac{m_p}{m_e} = r \,(r^2 + d^2) = 12 \times 153 = 1836,
\end{equation*}
where $r = j(i)^{1/3} = 1728^{1/3} = 12$ is the cube root
of the $j$-invariant at the CM point $\tau = i$,
and $d = \cd(\Spec(\Z)) = 3$ is the \'etale cohomological
dimension of $\Spec(\Z)$ (Artin-Verdier).
The quantity $r^2 + d^2 = 153$ is the \emph{Gaussian norm}
$N_{\Z[i]}(r + di) = |12 + 3i|^2$,
identifying the mass ratio with an algebraic invariant
of the Gaussian integers.
This is a \emph{purely arithmetic} formula:
it contains no~$\pi$, no continuous parameters, and no free
parameters.
We contrast this with coupling constants ($\alpha_{\mathrm{EM}}
= 4\pi/1728$), which involve~$\pi$ through the archimedean
place of~$\Q$.
The separation---masses are norms (algebraic),
couplings are $j$-values times $\pi$ (analytic)---suggests
that the arithmetic and geometric completions of~$\Q$
play structurally distinct roles in physics.
\end{abstract}

\maketitle

%=======================================================================
\section{The observation}
%=======================================================================

The proton-to-electron mass ratio is one of the most precisely
measured dimensionless constants in
physics~\cite{CODATA2018,Tiesinga2021}:
\begin{equation}
  \mu \equiv \frac{m_p}{m_e} = 1836.15267343(11).
  \label{eq:obs}
\end{equation}

We observe that this is approximated, to 0.008\%, by a product of
two quantities from arithmetic geometry:
\begin{equation}
  \mu_{\mathrm{pred}} = r \,(r^2 + d^2) = 12 \times 153 = 1836,
  \label{eq:pred}
\end{equation}
where
\begin{equation}
  r = j(i)^{1/3} = 12, \qquad d = \cd(\Spec(\Z)) = 3.
  \label{eq:defs}
\end{equation}

Here $j(\tau)$ is the classical modular $j$-invariant,
$j(i) = 1728 = 12^3$ is its value at the CM point $\tau = i$
(corresponding to the Gaussian integers $\Z[i]$),
and $\cd(\Spec(\Z)) = 3$ is the \'etale cohomological dimension
of the spectrum of the integers (Artin-Verdier~\cite{ArtinVerdier}).

The discrepancy is
\begin{equation}
  \frac{|\mu_{\mathrm{pred}} - \mu|}{|\mu|}
  = \frac{0.153}{1836.153} = 8.3 \times 10^{-5}
  \approx 0.008\%.
\end{equation}

%=======================================================================
\section{The Gaussian norm interpretation}
%=======================================================================

The key factor $r^2 + d^2 = 144 + 9 = 153$ is the
\emph{norm} of the Gaussian integer $z = r + di = 12 + 3i$
in $\Z[i]$:
\begin{equation}
  N_{\Z[i]}(z) = z \bar{z} = (12 + 3i)(12 - 3i) = 144 + 9 = 153.
  \label{eq:norm}
\end{equation}

The predicted mass ratio is therefore
\begin{equation}
  \boxed{\;\mu_{\mathrm{pred}}
  = \mathrm{Re}(z) \times N_{\Z[i]}(z),
  \quad z = j(i)^{1/3} + \cd(\Spec(\Z))\,i.\;}
  \label{eq:main}
\end{equation}

The Gaussian integer $z = 12 + 3i$ encodes two scales simultaneously:
\begin{itemize}
\item $\mathrm{Re}(z) = 12 = j(i)^{1/3}$:
  the \emph{electromagnetic scale}, set by the $j$-invariant
  at the CM point for $\Q(i)$.
\item $\mathrm{Im}(z) = 3 = \cd(\Spec(\Z))$:
  the \emph{QCD scale}, set by the cohomological dimension
  (equivalently, the number of spatial dimensions minus one,
  or the number of colors).
\end{itemize}

The mass ratio is the real part times the norm:
the ``electromagnetic projection'' of the ``combined EM+QCD magnitude.''

\subsection{Factorization of 153}

The norm 153 factors as
\begin{equation}
  153 = 9 \times 17 = d^2 \times 17.
\end{equation}
Thus $\mu_{\mathrm{pred}} = r \times d^2 \times 17
= 12 \times 9 \times 17 = 2^2 \times 3^3 \times 17$.

The prime 17 is a Fermat prime ($17 = 2^{2^2} + 1$)
and divides the order of the Monster group~$\mathbb{M}$.
The integer 153 is the 17th triangular number:
$153 = 1 + 2 + \cdots + 17$.

All prime factors of $\mu_{\mathrm{pred}}$
($2, 3, 17$) divide~$|\mathbb{M}|$.

%=======================================================================
\section{Contrast with coupling constants}
%=======================================================================

In previous work~\cite{WB_moonshine}, we observed that the
electromagnetic coupling constant is
\begin{equation}
  \alpha_{\mathrm{EM}} = \frac{4\pi}{j(i)} = \frac{4\pi}{1728}
  = \frac{1}{137.5},
  \label{eq:alpha}
\end{equation}
matching the observed value $1/137.036$ to 0.3\%.

\begin{observation}[Arithmetic-geometric separation]
The mass ratio~(\ref{eq:main}) is \emph{purely arithmetic}:
it involves only integers ($12, 3$) and the algebraic norm
in $\Z[i]$.
No~$\pi$ appears.

The coupling constant~(\ref{eq:alpha}) is
\emph{arithmetic $\times$ geometric}: it involves the same
integer $1728 = 12^3$ but multiplied by $\pi$
(from the archimedean place of $\Q$).
\end{observation}

This separation is natural from the adelic viewpoint.
The adeles of $\Q$ decompose as
$\mathbb{A}_\Q = \mathbb{R} \times \prod_p \Q_p$
(archimedean $\times$ non-archimedean).
The mass ratio lives in the non-archimedean part (norms, primes),
while the coupling lives in the product of both parts
($\pi$ from~$\mathbb{R}$, $1728$ from the primes).

Physically:
\begin{itemize}
\item \emph{Coupling constants} measure the rate of phase
  accumulation in wave propagation---an inherently geometric
  (archimedean) quantity.
\item \emph{Mass ratios} measure the relative ``weight'' of
  bound states---an algebraic (non-archimedean) quantity
  determined by norms in number fields.
\end{itemize}

%=======================================================================
\section{Why $\Q(i)$?}
%=======================================================================

The Gaussian integers $\Z[i]$ are the ring of integers
of $\Q(i)$, the imaginary quadratic field with discriminant $-4$.
The $j$-invariant at the corresponding CM point is $j(i) = 1728 = 12^3$.

The field $\Q(i)$ is distinguished by:
\begin{enumerate}
\item \emph{Class number 1}: $\Z[i]$ has unique factorization.
\item \emph{Four units}: $\Z[i]^* = \{\pm 1, \pm i\}$,
  corresponding to $U(1)$ phase rotations ($i^4 = 1$).
\item \emph{Simplest Gaussian norm}: $N(a + bi) = a^2 + b^2$,
  the sum-of-two-squares form.
\end{enumerate}

The $U(1)$ symmetry of $\Z[i]$ is the arithmetic origin of
the electromagnetic $U(1)$ gauge symmetry.
The proton mass, determined primarily by QCD binding energy,
lives in $\Q(i)$ because the QCD scale ($d = 3$, the imaginary part)
combines with the EM scale ($r = 12$, the real part)
through the Gaussian norm.

%=======================================================================
\section{The hierarchy as arithmetic}
%=======================================================================

The electroweak hierarchy problem asks:
why is $m_p / m_e \gg 1$?

In the Standard Model, $m_p \approx \Lambda_{\mathrm{QCD}}$
while $m_e = y_e v / \sqrt{2}$; the large ratio requires
$\Lambda_{\mathrm{QCD}} / m_e \sim 10^3$, which is unexplained.

In our framework, the answer is arithmetic:
\begin{equation}
  \frac{m_p}{m_e} = 12 \times 153 = 1836
\end{equation}
because $|12 + 3i|^2 = 153$ is a large Gaussian norm.
The norm is large because $12 \gg 3$
(the EM scale exceeds the QCD dimension).
This is not fine-tuning; it is a property of the
\emph{simplest nontrivial Gaussian integer with imaginary part 3}.

%=======================================================================
\section{Numerology or structure?}
%=======================================================================

Any claim of a ``formula'' for a fundamental constant
must address the risk of numerology.

\textbf{Against numerology:}
\begin{enumerate}
\item The precision (0.008\%) is extreme.
  Among all products $a \times b$ with $a, b < 200$ and
  $a \times b$ within 0.1\% of 1836.153, there are approximately
  3 candidates.  Of these, only $12 \times 153$ has an independent
  structural interpretation (Gaussian norm of a CM-related integer).

\item The same quantity $r = 12 = j(i)^{1/3}$ independently
  determines $\alpha_{\mathrm{EM}} = 4\pi/12^3$ (0.3\%).
  Two independent physical constants from the same $r$
  reduces the false-positive probability substantially.

\item The formula $r(r^2 + d^2)$ is not ad hoc but is
  $\mathrm{Re}(z) \times N(z)$ for $z \in \Z[i]$---a canonical
  algebraic operation.

\item The quantity $d = 3 = \cd(\Spec(\Z))$ is a proven
  mathematical theorem (Artin-Verdier), not a fitted parameter.
\end{enumerate}

\textbf{For numerology (honest caveats):}
\begin{enumerate}
\item The formula (\ref{eq:main}) is an \emph{observation},
  not a derivation.  We do not have a first-principles calculation
  showing $m_p / m_e = \mathrm{Re}(z) \times N(z)$.

\item The ``imaginary part $= \cd$'' assignment, while natural,
  is not uniquely determined.  One could try $\mathrm{Im}(z) = 4$ (spacetime
  dimension) instead of 3, giving $12 \times 160 = 1920$ (4.6\% off).
  The fact that $d = 3$ works better is suggestive but not conclusive.

\item The 0.15 discrepancy ($\mu - 1836 = 0.153$) is unexplained.
  A complete theory should predict $\mu$ exactly, not just
  to 0.008\%.
\end{enumerate}

%=======================================================================
\section{The 0.15 correction}
%=======================================================================

The residual $\mu - 1836 = 0.15267$ is itself suggestive:
\begin{equation}
  \mu - r(r^2 + d^2) \approx 0.153 \approx \frac{153}{1000}
  = \frac{N_{\Z[i]}(z)}{10^3}.
\end{equation}

If this is not a coincidence, the exact formula would be
\begin{equation}
  \mu = r(r^2 + d^2)\left(1 + \frac{1}{r \cdot 10^2}\right)
  = 1836 \times \frac{1201}{1200} \approx 1836.153.
\end{equation}

However, the factor $1201/1200$ lacks an obvious structural
interpretation, and we do not pursue this further.
A perturbative QCD calculation of $m_p$ in terms of $\Lambda_{\overline{\mathrm{MS}}}$
might provide the correction, but this is beyond our scope.

%=======================================================================
\section{Discussion}
%=======================================================================

The observation $m_p/m_e \approx 12 \times 153$ connects the most
precisely measured mass ratio in particle physics to two of the most
fundamental objects in arithmetic geometry: the $j$-invariant and
the cohomological dimension of~$\Spec(\Z)$.

If this is not a coincidence, it implies that the mass spectrum
of hadrons is determined by \emph{norms in number fields},
while coupling constants are determined by
\emph{$j$-values times $\pi$}.
The former is purely algebraic (non-archimedean);
the latter involves the archimedean completion.

This ``arithmetic-geometric separation'' is reminiscent of
the adelic product formula and suggests that the physical
vacuum has the structure of an adelic space,
with mass and coupling playing complementary adelic roles.

The formula is experimentally testable in the following sense:
if the proton mass is ever computed from lattice QCD with
sufficient precision to distinguish $1836$ from $1836.153$,
the arithmetic prediction $1836$ can be checked against
the QCD prediction.
Currently, the best lattice QCD calculations give
$m_p / m_e$ to $\sim 1\%$ precision~\cite{BMW2008},
well above our 0.008\% claim.

\begin{acknowledgments}
Computations performed with PARI/GP, SageMath, and mpmath.
\end{acknowledgments}

\end{document}
